% ====================================================================
% ANEXO: ARQUITECTURA TÉCNICA DETALLADA DE ChatHCE
% Trabajo Fin de Máster - AI Engineering
% Complemento técnico del Capítulo de Arquitectura
% ====================================================================
% Paquetes requeridos en el preámbulo del documento maestro:
%   inputenc(utf8), fontenc(T1), babel(spanish), amsmath, amssymb,
%   graphicx, booktabs, longtable, array, tabularx, enumitem, float,
%   listings, xcolor, hyperref, tikz, pgfplots
%   \definecolor{codebg}{gray}{0.95}  % <-- NECESARIO para listings
%   \usetikzlibrary{shapes.geometric,arrows.meta,positioning,calc,fit}
%   Requiere natbib o biblatex con comando \citep
% ====================================================================
%
% Paquetes necesarios: mismos que el capítulo principal
% ====================================================================

\chapter{Arquitectura técnica detallada de ChatHCE}
\label{appendix:arquitectura}

Este anexo proporciona la documentación técnica exhaustiva de todos los componentes del sistema ChatHCE, complementando la visión de alto nivel presentada en el Capítulo~\ref{cap:arquitectura}. Se incluyen descripciones detalladas de las herramientas, el pipeline \ac{RAG} completo, las capas de seguridad, los esquemas de base de datos, la gestión de modelos \ac{LLM}, el motor de visualización, y los servicios de infraestructura.

% ====================================================================
\section{Herramientas del agente: descripción detallada}
\label{sec:appendix_herramientas}
% ====================================================================

El agente unificado dispone de tres herramientas especializadas, cada una implementada como una subclase de \texttt{ClaudeToolAdapter} y registrada en el \texttt{ClaudeToolRegistry}. Esta sección documenta cada herramienta con sus parámetros, lógica de validación y ejemplos de uso.

\subsection{Interfaz ClaudeToolAdapter}

Todas las herramientas implementan la interfaz abstracta \texttt{ClaudeToolAdapter}, que define cuatro métodos obligatorios:

\begin{lstlisting}[language=Python, basicstyle=\tiny\ttfamily, backgroundcolor=\color{codebg}]
class ClaudeToolAdapter(ABC):
    @property
    @abstractmethod
    def name(self) -> str: ...
    
    @property
    @abstractmethod
    def description(self) -> str: ...
    
    @property
    @abstractmethod
    def args_schema(self) -> Type[BaseModel]: ...
    
    @abstractmethod
    def execute(self, **kwargs) -> Dict[str, Any]: ...
    
    def get_langchain_tool(self) -> StructuredTool:
        return StructuredTool.from_function(
            func=self.execute,
            name=self.name,
            description=self.description,
            args_schema=self.args_schema
        )
\end{lstlisting}

Este patrón \textit{Adapter} permite que cada herramienta defina su esquema Pydantic de parámetros (\texttt{args\_schema}), que Claude utiliza directamente para generar los argumentos \ac{JSON} de invocación.

\subsection{Database Tool (\texttt{query\_mimic\_database})}
\label{sec:appendix_dbtool}

La herramienta de base de datos proporciona acceso seguro y de solo lectura al dataset \ac{MIMIC}-IV-\ac{ED}. Implementa validación multicapa de consultas \ac{SQL} y soporta cinco tipos de consulta predefinidos.

\subsubsection{Tipos de Consulta}

\begin{table}[H]
\centering
\caption{Tipos de consulta soportados por la Database Tool con parámetros y límites.}
\label{tab:appendix_querytypes}
\footnotesize
\begin{tabularx}{\columnwidth}{llXl}
\toprule
\textbf{Tipo} & \textbf{Parámetro} & \textbf{Descripción} & \textbf{Límite} \\
\midrule
\texttt{patient\_summary} & \texttt{subject\_id} & Resumen completo: estancias, triaje, diagnósticos, signos vitales, medicamentos & 1000 filas \\
\addlinespace
\texttt{vital\_signs} & \texttt{stay\_id} & Signos vitales seriados de una estancia específica & 1000 filas \\
\addlinespace
\texttt{diagnoses} & \texttt{icd\_code} o \texttt{icd\_title} & Búsqueda de diagnósticos por código \ac{ICD} o título & 1000 filas \\
\addlinespace
\texttt{medications} & \texttt{subject\_id} & Historial de medicación habitual y dispensada en urgencias & 1000 filas \\
\addlinespace
\texttt{custom} & \texttt{custom\_query} & Consulta \ac{SQL} SELECT personalizada (validada) & 5000 filas \\
\bottomrule
\end{tabularx}
\end{table}

\subsubsection{Reglas de Validación SQL}

\begin{table}[H]
\centering
\caption{Reglas de validación de seguridad \ac{SQL} implementadas en la Database Tool.}
\label{tab:appendix_sqlsec}
\footnotesize
\begin{tabularx}{\columnwidth}{lX}
\toprule
\textbf{Regla} & \textbf{Descripción} \\
\midrule
Solo SELECT & La consulta debe comenzar con SELECT; se bloquean INSERT, UPDATE, DELETE, DROP, ALTER, CREATE, TRUNCATE, GRANT, REVOKE \\
Límite de longitud & Máximo 2000 caracteres por consulta \\
Sin comentarios & Se bloquean \texttt{-{}-} y \texttt{/* */} para prevenir inyección por comentarios \\
Sin multi-sentencia & Se bloquea \texttt{;} para prevenir ejecución de múltiples sentencias \\
Tautologías & Detección de patrones: \texttt{OR 1=1}, \texttt{OR 'x'='x'}, \texttt{OR TRUE}, \texttt{OR 1}, \texttt{OR (1=1)}, \texttt{OR 2=2}, \texttt{OR 1<>0} \\
Tablas permitidas & Solo: \texttt{edstays}, \texttt{triage}, \texttt{vitalsign}, \texttt{diagnosis}, \texttt{medrecon}, \texttt{pyxis} \\
Funciones peligrosas & Bloqueo de: \texttt{PG\_SLEEP}, \texttt{PG\_READ\_FILE}, \texttt{DBLINK}, \texttt{LO\_IMPORT}, \texttt{EXEC}, \texttt{SYSTEM} \\
Tablas del sistema & Bloqueo de acceso a \texttt{pg\_*} e \texttt{information\_schema} \\
Complejidad máxima & Máximo 10 condiciones (WHERE + AND + OR) \\
Columnas erróneas & Detección y corrección: \texttt{drugname}$\to$\texttt{name}, \texttt{patient\_id}$\to$\texttt{subject\_id} \\
\bottomrule
\end{tabularx}
\end{table}

\subsubsection{Implementación de Validación}

\begin{lstlisting}[language=Python, basicstyle=\tiny\ttfamily, backgroundcolor=\color{codebg}]
BLACKLISTED_OPERATIONS = {
    'INSERT', 'UPDATE', 'DELETE', 'DROP', 'CREATE', 'ALTER',
    'TRUNCATE', 'GRANT', 'REVOKE', 'EXECUTE', 'COPY'
}

TAUTOLOGY_PATTERNS = [
    r"OR\s+1\s*=\s*1",
    r"OR\s+'[^']*'\s*=\s*'[^']*'",
    r"OR\s+\d+\s*=\s*\d+",
]

def validate_query(sql: str) -> bool:
    sql_upper = sql.upper().strip()
    if not sql_upper.startswith('SELECT'):
        raise SecurityError("Solo se permiten consultas SELECT")
    for op in BLACKLISTED_OPERATIONS:
        if op in sql_upper:
            raise SecurityError(f"Operacion {op} no permitida")
    for pattern in TAUTOLOGY_PATTERNS:
        if re.search(pattern, sql, re.IGNORECASE):
            raise SecurityError("Patron de inyeccion detectado")
    return True
\end{lstlisting}

\subsection{RAG Tool (\texttt{search\_clinical\_documents})}
\label{sec:appendix_ragtool}

La herramienta \ac{RAG} proporciona búsqueda semántica en documentos clínicos indexados. Acepta tres parámetros: \texttt{query} (obligatorio, la consulta en lenguaje natural), \texttt{specialty} (opcional, filtro por especialidad médica) y \texttt{top\_k} (opcional, por defecto 8).

\subsubsection{Pipeline Interno}

El flujo de ejecución del \ac{RAG} Tool sigue siete etapas secuenciales:

\begin{enumerate}[nosep, leftmargin=*]
  \item \textbf{Recepción}: Validación de parámetros con Pydantic.
  \item \textbf{Augmentación}: \texttt{QueryAugmenter} genera 3 consultas alternativas con Claude Haiku (Multi-Query, temperatura 0.4) y un documento hipotético (\ac{HyDE}, temperatura 0.2).
  \item \textbf{Búsqueda}: \texttt{ImprovedRAGService.search()} ejecuta búsqueda híbrida para cada consulta augmentada (\texttt{fetch\_k = top\_k $\times$ 4 = 32}).
  \item \textbf{Deduplicación}: Eliminación de duplicados por hash de contenido.
  \item \textbf{Reranking}: Cross-encoder \texttt{ms-marco-MiniLM-L-6-v2} reordena los candidatos; selecciona los top-8.
  \item \textbf{Expansión}: Recuperación de chunks padre (3000 chars) mediante \texttt{parent\_id} para contexto completo.
  \item \textbf{Formateo}: Generación de salida con citación de fuentes (nombre de documento, sección, score).
\end{enumerate}

\subsubsection{Documentos Clínicos Indexados}

\begin{table}[H]
\centering
\caption{Documentos clínicos indexados en el sistema \ac{RAG} con descripción detallada.}
\label{tab:appendix_ragdocs}
\footnotesize
\begin{tabularx}{\columnwidth}{lX}
\toprule
\textbf{Documento} & \textbf{Descripción} \\
\midrule
Manual del Residente de Medicina Intensiva & Hospital Universitario Virgen del Rocío. Formación MIR, competencias, rotaciones, habilidades técnicas en Medicina Intensiva. \\
\addlinespace
Estándares y Recomendaciones para \ac{UCI} & Ministerio de Sanidad de España. Ratios médico-paciente, niveles asistenciales, indicadores de calidad, seguridad del paciente en \ac{UCI}. \\
\addlinespace
El Libro de la \ac{UCI} & Paul L. Marino, 3.ª edición. Referencia clínica completa de cuidados críticos: ventilación mecánica, sepsis, hemostasia, nutrición en \ac{UCI}. \\
\bottomrule
\end{tabularx}
\end{table}

\subsubsection{Formato de Salida}

El \ac{RAG} Tool formatea la salida para Claude con bloques delimitados por documento:

\begin{lstlisting}[basicstyle=\tiny\ttfamily, backgroundcolor=\color{codebg}]
--- Documento 1: Manual del Residente UCI ---
Fuente: manual_residente_uci.pdf (pag. 45)
Score: 0.847
Contenido: [hasta 1500 caracteres del chunk padre]

--- Documento 2: Estandares UCI ---
Fuente: estandares_uci_ministerio.pdf (pag. 23)
Score: 0.792
Contenido: [hasta 1500 caracteres del chunk padre]
\end{lstlisting}

El campo \texttt{raw\_observation} preserva el texto completo sin truncar, permitiendo que el evaluador \ac{RAGAS} parsee directamente los contextos recuperados.

\subsection{Visualization Tool (\texttt{request\_visualization})}
\label{sec:appendix_viztool}

La herramienta de visualización genera gráficas interactivas con Plotly. Acepta parámetros para el tipo de visualización (\texttt{viz\_type}), la consulta del usuario (\texttt{query}), y opcionalmente datos precargados.

\subsubsection{Catálogo Completo de Plantillas}

\begin{table}[H]
\centering
\caption{Catálogo completo de plantillas de visualización predefinidas.}
\label{tab:appendix_templates}
\footnotesize
\begin{tabularx}{\columnwidth}{lX}
\toprule
\textbf{Plantilla} & \textbf{Caso de uso clínico} \\
\midrule
\texttt{vitals\_timeline} & Evolución temporal de signos vitales de un paciente \\
\texttt{vitals\_comparison} & Comparación de múltiples signos vitales simultáneos \\
\texttt{distribution\_histogram} & Distribución de valores (edad, frecuencia cardíaca, SpO2) \\
\texttt{acuity\_pie} & Distribución de niveles de acuidad \ac{ESI} \\
\texttt{medication\_bar} & Frecuencia de medicamentos administrados \\
\texttt{diagnosis\_treemap} & Jerarquía de diagnósticos \ac{ICD} \\
\texttt{correlation\_heatmap} & Correlación entre variables clínicas \\
\texttt{boxplot\_comparison} & Comparación de distribuciones por grupo \\
\texttt{scatter\_vitals} & Relación entre pares de signos vitales \\
\texttt{los\_distribution} & Distribución de duración de estancia \\
\texttt{triage\_radar} & Perfil multidimensional de triaje \\
\texttt{timeline\_events} & Línea temporal de eventos clínicos \\
\bottomrule
\end{tabularx}
\end{table}

\subsubsection{Lógica de Selección Automática}

\begin{lstlisting}[language=Python, basicstyle=\tiny\ttfamily, backgroundcolor=\color{codebg}]
def select_template(data: pd.DataFrame, viz_type: str) -> str:
    has_time = any(col for col in data.columns 
                   if 'time' in col.lower() or 'date' in col.lower())
    vital_cols = ['heartrate', 'temperature', 'resprate', 
                  'o2sat', 'sbp', 'dbp']
    has_vitals = any(col in data.columns for col in vital_cols)
    
    if has_time and has_vitals:
        return 'vitals_timeline'
    elif viz_type == 'distribution':
        return 'distribution_histogram'
    elif viz_type == 'comparison':
        return 'vitals_comparison'
    # ... logica adicional para otros tipos
    else:
        return None  # Fallback a LLM
\end{lstlisting}

\subsubsection{Palabras Clave de Activación}

\begin{table}[H]
\centering
\caption{Palabras clave que activan automáticamente la herramienta de visualización.}
\label{tab:appendix_vizkeywords}
\footnotesize
\begin{tabularx}{\columnwidth}{lX}
\toprule
\textbf{Categoría} & \textbf{Palabras clave} \\
\midrule
Visualización explícita & gráfico, gráfica, visualiza, visualización, muestra un gráfico, genera una gráfica, dibuja, plot, chart, graph \\
Consultas temporales & evolución, tendencia, progreso, cambio, a lo largo del tiempo, durante, histórico, timeline \\
Distribuciones & distribución de, frecuencia de, histograma, boxplot \\
Estadísticas & estadísticas de, análisis de, comparación de \\
\bottomrule
\end{tabularx}
\end{table}

\subsubsection{Fallback a Generación LLM}

Cuando no hay plantilla disponible (aproximadamente 10\% de los casos), el sistema escala a Claude Sonnet 4.5 para generar código Plotly personalizado. El código generado se valida mediante parsing \ac{AST} antes de ejecutarse en un sandbox aislado con timeout de 10 segundos.

\begin{lstlisting}[language=Python, basicstyle=\tiny\ttfamily, backgroundcolor=\color{codebg}]
def generate_with_llm(data: pd.DataFrame, request: str) -> str:
    prompt = f"""
    Genera codigo Python con Plotly para visualizar los siguientes datos.
    Columnas disponibles: {list(data.columns)}
    Primeras filas: {data.head(3).to_dict()}
    Solicitud del usuario: {request}
    El codigo debe:
    1. Usar solo plotly.express o plotly.graph_objects
    2. Retornar un objeto fig de tipo Figure
    3. No usar plt.show() ni fig.show()
    4. Incluir titulos y etiquetas en espanol
    """
    return self.llm.generate(prompt, model='claude-sonnet-4-5')
\end{lstlisting}

% ====================================================================
\section{Pipeline \ac{RAG}: detalles técnicos completos}
\label{sec:appendix_rag}
% ====================================================================

Esta sección complementa la descripción del pipeline \ac{RAG} presentada en la Sección~\ref{sec:rag} del capítulo principal, proporcionando todos los detalles técnicos de implementación.

\subsection{Configuración del chunking y embeddings}

\begin{table}[H]
\centering
\caption{Configuración completa del sistema de chunking y embeddings.}
\label{tab:appendix_chunking}
\footnotesize
\begin{tabularx}{\columnwidth}{llX}
\toprule
\textbf{Parámetro} & \textbf{Valor} & \textbf{Descripción} \\
\midrule
\texttt{parent\_chunk\_size} & 3000 chars & Tamaño del chunk padre para contexto completo \\
\texttt{child\_chunk\_size} & 800 chars & Tamaño del chunk hijo para búsqueda precisa \\
Modelo de embeddings & \texttt{all-MiniLM-L6-v2} & sentence-transformers, 384 dimensiones \\
Dispositivo & CUDA (fallback CPU) & Detección automática de GPU disponible \\
Normalización & True & Embeddings normalizados para similitud coseno \\
Singleton & Sí & Una sola instancia por proceso Python (evita recarga CUDA) \\
\texttt{top\_k} búsqueda & 8 (con reranking) & Número de resultados finales \\
\texttt{fetch\_k} & top\_k $\times$ 4 = 32 & Candidatos evaluados por el reranker \\
Truncado de contenido & 1500 chars & Máximo de caracteres pasados al \ac{LLM} por documento \\
\bottomrule
\end{tabularx}
\end{table}

\subsection{Búsqueda híbrida: implementación RPC en PostgreSQL}

La búsqueda híbrida se ejecuta mediante una función RPC de PostgreSQL que combina las dos modalidades en una sola query:

\begin{itemize}[nosep, leftmargin=*]
  \item \textbf{Búsqueda vectorial} (pgvector): Similitud coseno entre el embedding de la consulta (384d) y los embeddings de los chunks hijo. Utiliza índice \ac{HNSW} para búsqueda ANN en tiempo sublineal.
  \item \textbf{Búsqueda full-text} (tsvector): Búsqueda léxica con stemming en español (\texttt{spanish} configuration). Captura coincidencias exactas de términos médicos, códigos \ac{ICD}, nombres de fármacos.
\end{itemize}

La fusión \ac{RRF} combina los rankings sin requerir calibración de pesos:

\begin{equation}
\text{RRF}(d) = \sum_{r \in R} \frac{1}{k + \text{rank}_r(d)}
\end{equation}

donde $R = \{\text{vectorial}, \text{léxica}\}$ y $k = 60$ es la constante de suavizado que balancea la influencia de las primeras posiciones. El valor $k=60$ es el estándar recomendado en la literatura \citep{cormack2009rrf} y fue adoptado sin necesidad de ajuste adicional.

\subsection{Reranking por cross-encoder}

El reranker utiliza el modelo \texttt{cross-encoder/ms-marco-MiniLM-L-6-v2}. A diferencia de los bi-encoders (que generan embeddings independientes), el cross-encoder procesa la consulta y el documento conjuntamente, produciendo una puntuación de relevancia más precisa. Si el modelo no está disponible, el sistema continúa sin reranking de forma transparente.

\subsection{Augmentación de consultas: detalles}

\textbf{Multi-Query Generation.} Genera hasta 3 consultas alternativas desde diferentes perspectivas médicas (diagnóstico, tratamiento, fisiopatología), utilizando terminología variada y sinónimos clínicos. Temperatura: 0.4 para favorecer diversidad léxica.

\textbf{HyDE (Hypothetical Document Embeddings).} Genera un fragmento hipotético de guía clínica que respondería a la consulta, con temperatura 0.2 para mayor precisión factual. El embedding de este fragmento hipotético captura mejor la semántica del resultado esperado que el embedding de la pregunta original.

\subsection{Diagnóstico y mejoras del pipeline RAG}
\label{sec:appendix_rag_mejoras}

Durante la evaluación con el framework \ac{RAGAS}, se identificaron seis causas raíz que impedían el correcto funcionamiento del pipeline:

\subsubsection{Causa 1: Desalineación del System Prompt}

El system prompt definía cuándo usar \texttt{search\_clinical\_documents} con términos de ``urgencias y emergencias'', pero los documentos indexados versan sobre \textbf{UCI y Medicina Intensiva}. El agente no reconocía que debía buscar en los documentos para preguntas como ``¿Cuántas camas tiene la \ac{UCI}?''.

\textbf{Solución:} Se reescribió la sección de selección de herramientas para incluir explícitamente \ac{UCI}, cuidados críticos y medicina intensiva, listar los documentos reales indexados, y añadir 8 ejemplos concretos de preguntas que deben usar la herramienta. La directiva pasó de ``considera usar'' a ``SIEMPRE invoca''.

\subsubsection{Causa 2: Límite de Tokens Insuficiente}

El \texttt{PromptManager} se inicializaba con \texttt{max\_tokens=4000}. El prompt completo (schema de BD + directivas + tool descriptions) alcanzaba 7726 tokens, truncándose y eliminando parte del schema y las directivas de selección de herramientas.

\textbf{Solución:} Incremento de \texttt{max\_tokens} a 8000, eliminando el truncado del prompt.

\subsubsection{Causa 3: Contenido Truncado a 600 Caracteres}

El método \texttt{format\_output} del \ac{RAG} Tool truncaba el contenido de cada documento a 600 caracteres. Los parent chunks ($\sim$3000 chars) quedaban cortados, impidiendo que el \ac{LLM} encontrara datos específicos.

\textbf{Solución:} Incremento del límite de truncado a 1500 caracteres.

\subsubsection{Causa 4: top\_k Insuficiente}

Con \texttt{top\_k=5} y query augmentation, los 5 resultados finales eran insuficientes para cubrir fragmentos muy específicos dispersos en documentos de cientos de páginas.

\textbf{Solución:} Incremento de \texttt{top\_k} a 8.

\subsubsection{Causa 5: fetch\_k Bajo en el Reranker}

\texttt{fetch\_k = top\_k $\times$ 3} significaba que con \texttt{top\_k=5} solo se evaluaban 15 candidatos antes del reranking.

\textbf{Solución:} Incremento del factor a 4, resultando en \texttt{fetch\_k = top\_k $\times$ 4 = 32} candidatos.

\subsubsection{Causa 6: Contexto \ac{RAG} No Llegaba a RAGAS}

El campo \texttt{raw\_output} se truncaba a 1000 caracteres en \texttt{\_extract\_tool\_result}, y el evaluador \ac{RAGAS} intentaba parsear un campo \texttt{raw\_observation} que no existía.

\textbf{Solución:} Preservación del texto completo en \texttt{raw\_output} y adición del campo \texttt{raw\_observation} para el evaluador.

\subsection{Evolución completa de métricas RAGAS}

\begin{table}[H]
\centering
\caption{Evolución de métricas \ac{RAGAS} en el golden set DB (40 preguntas sobre \ac{MIMIC}-IV-\ac{ED}).}
\label{tab:appendix_ragas_db}
\footnotesize
\begin{tabularx}{\columnwidth}{lrrr}
\toprule
\textbf{Métrica} & \textbf{Antes (27-abr)} & \textbf{Después (28-abr)} & \textbf{$\Delta$} \\
\midrule
Faithfulness & 0.7995 & 0.7357 & $-$0.064 \\
Answer Relevancy & 0.5011 & 0.5944 & \textbf{+0.093} \\
Context Precision & 0.7417 & 0.6021 & $-$0.140 \\
Context Recall & 0.7917 & \textbf{0.8458} & \textbf{+0.054} \\
\bottomrule
\end{tabularx}
\end{table}

\begin{table}[H]
\centering
\caption{Evolución de métricas \ac{RAGAS} en el golden set \ac{RAG} (30 preguntas sobre guías clínicas \ac{UCI}).}
\label{tab:appendix_ragas_rag}
\footnotesize
\begin{tabularx}{\columnwidth}{lrrr}
\toprule
\textbf{Métrica} & \textbf{Sin docs} & \textbf{Con docs sin mejoras} & \textbf{Con mejoras} \\
\midrule
Faithfulness & 0.717 & 0.788 & 0.571 \\
Answer Relevancy & 0.000 & 0.096 & 0.435 \\
Context Precision & 0.017 & 0.011 & \textbf{0.450} \\
Context Recall & 0.033 & 0.067 & \textbf{0.469} \\
\bottomrule
\end{tabularx}
\end{table}

\subsection{Análisis de métricas pendientes}

\textbf{Faithfulness (\ac{RAG}: 0.57).} Las respuestas incluyen interpretaciones clínicas del \ac{LLM} que van más allá del texto recuperado. Posibles mejoras: reforzar directivas anti-alucinación específicas para respuestas \ac{RAG}; añadir instrucción ``responde SOLO con lo que aparece en los fragmentos''.

\textbf{Answer Relevancy (\ac{RAG}: 0.48).} Las respuestas son extensas y añaden contexto no solicitado. Posibles mejoras: directiva de concisión específica; limitar longitud cuando la pregunta es directa.

\textbf{Context Precision (\ac{RAG}: 0.55).} Aproximadamente el 55\% del contexto recuperado es directamente relevante. Posibles mejoras: ajustar el modelo de reranking para el dominio médico en español; aumentar peso de búsqueda léxica para términos técnicos exactos.

\textbf{Context Recall (\ac{RAG}: 0.49).} Aproximadamente la mitad de la información necesaria se recupera. Posibles mejoras: aumentar \texttt{top\_k} a 10-12 para preguntas directas; mejorar chunking para preservar datos numéricos en contexto.

% ====================================================================
\section{Seguridad: detalles de implementación}
\label{sec:appendix_seguridad}
% ====================================================================

Esta sección detalla la implementación de las seis capas de seguridad presentadas en la Sección~\ref{sec:seguridad} del capítulo principal.

\subsection{Capa 1: rate limiting}

El \texttt{RateLimiter} implementa un algoritmo de ventana deslizante con tres niveles de protección y bloqueo progresivo ante violaciones repetidas.

\begin{table}[H]
\centering
\caption{Configuración completa del Rate Limiter con duraciones de bloqueo progresivo.}
\label{tab:appendix_ratelimit}
\footnotesize
\begin{tabularx}{\columnwidth}{llX}
\toprule
\textbf{Límite} & \textbf{Valor} & \textbf{Descripción} \\
\midrule
Burst & 5 requests / 10s & Protección contra ráfagas \\
Por minuto & 30 requests / min & Límite de uso normal \\
Por hora & 300 requests / hora & Límite de uso sostenido \\
Longitud máxima & 5000 chars & Longitud máxima del mensaje \\
Bloqueo 1.ª violación & 30s & Primer bloqueo \\
Bloqueo 2.ª violación & 60s & Segundo bloqueo \\
Bloqueo 3.ª violación & 120s & Tercer bloqueo \\
Bloqueo 4.ª violación & 300s & Cuarto bloqueo \\
Bloqueo 5.ª violación & 600s & Máximo bloqueo \\
Limpieza automática & 300s & Intervalo de limpieza de entradas expiradas \\
\bottomrule
\end{tabularx}
\end{table}

\subsection{Capa 2: validación de entrada}

La primera línea de defensa opera sobre la consulta del usuario antes de cualquier procesamiento:

\begin{itemize}[nosep, leftmargin=*]
  \item \textbf{Longitud máxima}: 5000 caracteres para prevenir ataques de denegación de servicio.
  \item \textbf{Sanitización}: Eliminación de caracteres de control y secuencias potencialmente maliciosas.
  \item \textbf{Formato}: Verificación de codificación UTF-8 válida.
\end{itemize}

\subsection{Capa 3: validación \ac{SQL} multicapa}

Documentada en detalle en la Sección~\ref{sec:appendix_dbtool}. Incluye lista negra de operaciones DDL/DML, detección de tautologías (10 patrones), inyección por comentarios, restricción a tablas permitidas del esquema \texttt{mimic\_ed}, bloqueo de funciones peligrosas, y límite de complejidad.

\subsection{Capa 4: autenticación JWT y row level security}

Supabase Auth proporciona autenticación JWT con RLS habilitado en las 13 tablas del sistema (6 tablas \ac{MIMIC}-IV-\ac{ED} + 7 tablas de aplicación). Las políticas RLS garantizan que cada usuario solo puede acceder a sus propias sesiones, mensajes y preferencias.

\subsection{Capa 5: audit logging}

Todas las operaciones del sistema se registran con logs estructurados que incluyen: timestamp, session ID, tipo de operación, herramienta invocada, tiempo de ejecución, resultado (éxito/error), y modelo \ac{LLM} utilizado. Estos logs permiten trazabilidad completa para auditoría y debugging.

\subsection{Capa 6: directivas anti-alucinación}

\begin{table}[H]
\centering
\caption{Directivas anti-alucinación implementadas en el prompt del sistema.}
\label{tab:appendix_antihallucination}
\footnotesize
\begin{tabularx}{\columnwidth}{lX}
\toprule
\textbf{Directiva} & \textbf{Descripción} \\
\midrule
Prohibiciones absolutas & NUNCA inventar subject\_id, stay\_id, hadm\_id; NUNCA fabricar signos vitales, diagnósticos, medicamentos, fechas \\
Datos faltantes & Si no encuentra datos: ``No encontré información sobre [X]''; si hay nulos: ``Algunos registros tienen datos incompletos'' \\
Citación de fuentes & Siempre indicar herramienta usada y tabla de origen; para \ac{RAG}, citar documento fuente \\
Incertidumbre & Distinguir datos verificados de interpretaciones; usar ``Según los datos disponibles...'' \\
Limitaciones & \ac{MIMIC}-IV-\ac{ED} es dataset demo con 222 pacientes; solo urgencias; sin laboratorio, imagen ni cirugía \\
Concisión & Responder únicamente lo preguntado; no añadir información no solicitada \\
Fidelidad & Basar respuesta exclusivamente en datos de herramientas; no usar conocimiento general \\
\bottomrule
\end{tabularx}
\end{table}

\subsection{Sandbox \ac{AST} para visualizaciones}

El código Plotly generado por el \ac{LLM} se ejecuta en un entorno aislado con validación \ac{AST}:

\begin{lstlisting}[language=Python, basicstyle=\tiny\ttfamily, backgroundcolor=\color{codebg}]
ALLOWED_MODULES = {'plotly', 'pandas', 'numpy', 'math'}
FORBIDDEN_BUILTINS = {'exec', 'eval', 'compile', 'open', '__import__'}

def validate_code_ast(code: str) -> bool:
    tree = ast.parse(code)
    for node in ast.walk(tree):
        if isinstance(node, ast.Import):
            for alias in node.names:
                if alias.name.split('.')[0] not in ALLOWED_MODULES:
                    raise SecurityError(f"Modulo {alias.name} no permitido")
        if isinstance(node, ast.Call):
            if isinstance(node.func, ast.Name):
                if node.func.id in FORBIDDEN_BUILTINS:
                    raise SecurityError(f"Funcion {node.func.id} prohibida")
    return True
\end{lstlisting}

\subsection{Resultados de pruebas de seguridad}

\begin{table}[H]
\centering
\caption{Resultados detallados de las pruebas de seguridad (28 de abril de 2026).}
\label{tab:appendix_secresults}
\footnotesize
\begin{tabularx}{\columnwidth}{llllX}
\toprule
\textbf{Categoría} & \textbf{Tests} & \textbf{Pasados} & \textbf{Score} & \textbf{Resultado} \\
\midrule
\ac{SQL} Injection & 7 & 7 & 1.000 & Superado \\
Prompt Injection & 3 & 2 & 0.667 & No superado \\
Anti-Alucinación & 3 & 3 & 1.000 & Superado \\
\midrule
\textbf{Total} & \textbf{13} & \textbf{12} & \textbf{0.923} & \textbf{12/13 superados} \\
\bottomrule
\end{tabularx}
\end{table}

El test fallido (SEC-PROMPT-003) corresponde a una solicitud de generación de datos ficticios de pacientes. El agente rechazó correctamente la solicitud, pero el evaluador consideró que la respuesta no fue suficientemente explícita en la negativa.

% ====================================================================
\section{Arquitectura de base de datos: detalles completos}
\label{sec:appendix_db}
% ====================================================================

\subsection{Tablas \ac{MIMIC}-IV-ED}

\begin{table}[H]
\centering
\caption{Tabla \texttt{edstays}: columnas, tipos y descripción clínica completa.}
\label{tab:appendix_edstays}
\footnotesize
\begin{tabularx}{\columnwidth}{llX}
\toprule
\textbf{Columna} & \textbf{Tipo} & \textbf{Descripción} \\
\midrule
\texttt{subject\_id} & INT & Identificador único del paciente \\
\texttt{stay\_id} & INT (\ac{PK}) & Identificador único de la estancia \\
\texttt{hadm\_id} & NUMERIC & Identificador de admisión hospitalaria \\
\texttt{intime} & TIMESTAMP & Fecha/hora de entrada en urgencias \\
\texttt{outtime} & TIMESTAMP & Fecha/hora de salida de urgencias \\
\texttt{gender} & VARCHAR & Género del paciente (M/F) \\
\texttt{race} & VARCHAR & Raza/etnia del paciente \\
\texttt{arrival\_transport} & VARCHAR & Medio de transporte de llegada \\
\texttt{disposition} & VARCHAR & Destino al alta (home, admitted, expired, etc.) \\
\bottomrule
\end{tabularx}
\end{table}

\subsection{Tabla Vectorial \ac{RAG} (\texttt{rag\_chunks})}

\begin{table}[H]
\centering
\caption{Estructura completa de la tabla \texttt{rag\_chunks}.}
\label{tab:appendix_ragchunks}
\footnotesize
\begin{tabularx}{\columnwidth}{llX}
\toprule
\textbf{Columna} & \textbf{Tipo} & \textbf{Propósito} \\
\midrule
\texttt{id} & \texttt{uuid} & Clave primaria (\texttt{gen\_random\_uuid()}) \\
\texttt{chunk\_id} & \texttt{text} & Identificador único del chunk \\
\texttt{content} & \texttt{text} & Contenido textual (padre: 3000 chars, hijo: 800 chars) \\
\texttt{embedding} & \texttt{vector(384)} & Embedding \texttt{all-MiniLM-L6-v2}. Índice \ac{HNSW} \\
\texttt{fts} & \texttt{tsvector} & Columna generada desde \texttt{content} con config. \texttt{spanish}. Índice \ac{GIN} \\
\texttt{is\_parent} & \texttt{boolean} & \texttt{true} = chunk padre, \texttt{false} = chunk hijo \\
\texttt{parent\_id} & \texttt{text} & Referencia al chunk padre \\
\texttt{filename} & \texttt{text} & Nombre del documento fuente para citación \\
\texttt{document\_id} & \texttt{uuid} & Identificador del documento (eliminación en cascada) \\
\texttt{metadata} & \texttt{jsonb} & Metadatos: especialidad, tipo, página, fecha \\
\texttt{created\_at} & \texttt{timestamptz} & Marca temporal de indexación \\
\bottomrule
\end{tabularx}
\end{table}

\subsection{Estrategia de indexación}

\begin{table}[H]
\centering
\caption{Índices implementados en Supabase con tipo y propósito.}
\label{tab:appendix_indices}
\footnotesize
\begin{tabularx}{\columnwidth}{lllX}
\toprule
\textbf{Tabla} & \textbf{Columna(s)} & \textbf{Tipo} & \textbf{Propósito} \\
\midrule
\texttt{edstays} & \texttt{subject\_id} & B-tree & Consultas por paciente \\
\texttt{edstays} & \texttt{intime} & B-tree & Consultas temporales \\
\texttt{triage} & \texttt{subject\_id} & B-tree & Triaje por paciente \\
\texttt{triage} & \texttt{acuity} & B-tree & Filtrado por urgencia \\
\texttt{vitalsign} & \texttt{subject\_id} & B-tree & Signos vitales por paciente \\
\texttt{vitalsign} & \texttt{(subject\_id, charttime)} & B-tree comp. & Tendencias temporales \\
\texttt{diagnosis} & \texttt{subject\_id} & B-tree & Diagnósticos por paciente \\
\texttt{diagnosis} & \texttt{icd\_code} & B-tree & Búsqueda por código \ac{ICD} \\
\texttt{diagnosis} & \texttt{(subject\_id, stay\_id)} & B-tree comp. & Diagnósticos por estancia \\
\texttt{medrecon} & \texttt{(subject\_id, charttime)} & B-tree comp. & Medicación temporal \\
\texttt{pyxis} & \texttt{(subject\_id, charttime)} & B-tree comp. & Dispensación temporal \\
\texttt{chat\_messages} & \texttt{session\_id} & B-tree & Mensajes por sesión \\
\texttt{chat\_sessions} & \texttt{user\_id} & B-tree & Sesiones por usuario \\
\texttt{rag\_chunks} & \texttt{embedding} & \ac{HNSW} & Búsqueda ANN vectorial \\
\texttt{rag\_chunks} & \texttt{fts} & \ac{GIN} & Full-text en español \\
\bottomrule
\end{tabularx}
\end{table}

\subsection{Valores de referencia clínicos}

\begin{table}[H]
\centering
\caption{Valores de referencia clínicos normales implementados en el sistema.}
\label{tab:appendix_vitales}
\footnotesize
\begin{tabularx}{\columnwidth}{llll}
\toprule
\textbf{Parámetro} & \textbf{Columna} & \textbf{Rango normal} & \textbf{Unidad} \\
\midrule
Temperatura & \texttt{temperature} & 96.8 -- 100.4 & °F \\
Frecuencia cardíaca & \texttt{heartrate} & 60 -- 100 & lpm \\
Frecuencia respiratoria & \texttt{resprate} & 12 -- 20 & rpm \\
Saturación de oxígeno & \texttt{o2sat} & 95 -- 100 & \% \\
Presión arterial sistólica & \texttt{sbp} & 90 -- 140 & mmHg \\
Presión arterial diastólica & \texttt{dbp} & 60 -- 90 & mmHg \\
\bottomrule
\end{tabularx}
\end{table}

\begin{table}[H]
\centering
\caption{Niveles de acuidad del triaje (Emergency Severity Index).}
\label{tab:appendix_acuidad}
\footnotesize
\begin{tabularx}{\columnwidth}{clX}
\toprule
\textbf{Nivel} & \textbf{Categoría} & \textbf{Tiempo de atención} \\
\midrule
1 & Resucitación & Inmediato --- riesgo vital inminente \\
2 & Emergencia & 10--15 minutos \\
3 & Urgente & 30--60 minutos \\
4 & Menos urgente & 1--2 horas \\
5 & No urgente & 2--4 horas \\
\bottomrule
\end{tabularx}
\end{table}

% ====================================================================
\section{Gestor de modelos \ac{LLM}: implementación detallada}
\label{sec:appendix_llm}
% ====================================================================

\subsection{ClaudeLLMManager: configuración completa}

\begin{table}[H]
\centering
\caption{Configuración completa de la cadena de modelos \ac{LLM}.}
\label{tab:appendix_models}
\footnotesize
\begin{tabularx}{\columnwidth}{llrrrX}
\toprule
\textbf{Rol} & \textbf{Modelo} & \textbf{max\_tok} & \textbf{temp} & \textbf{timeout} & \textbf{Uso} \\
\midrule
Primario & \texttt{haiku-4-5-20251001} & 4096 & 0.1 & 30s & Agente principal \\
Secundario & \texttt{sonnet-4-5-20250929} & 8192 & 0.1 & 45s & Fallback \\
Terciario & \texttt{opus-4-20250514} & 4096 & 0.1 & 60s & Fallback final \\
Visualización & \texttt{sonnet-4-5} & 8192 & 0.1 & 45s & Código Plotly \\
\ac{RAG} & \texttt{haiku-4-5} & 4000 & 0.1 & 60s & Query augment. \\
\bottomrule
\end{tabularx}
\end{table}

\subsection{Lógica de reintentos y escalado}

La implementación utiliza la biblioteca \texttt{tenacity} para reintentos con backoff exponencial:

\begin{lstlisting}[language=Python, basicstyle=\tiny\ttfamily, backgroundcolor=\color{codebg}]
@retry(
    stop=stop_after_attempt(3),
    wait=wait_exponential(multiplier=2.0, min=1, max=60),
    retry=retry_if_exception_type((RateLimitError, APIError))
)
def _call_with_fallback(self, messages: List[Dict]) -> str:
    for model in self.model_chain:
        try:
            return self._call_model(model, messages)
        except ModelOverloadedError:
            continue
    raise AllModelsExhaustedError()
\end{lstlisting}

\begin{table}[H]
\centering
\caption{Parámetros de configuración de reintentos.}
\label{tab:appendix_reintentos}
\footnotesize
\begin{tabularx}{\columnwidth}{Xr}
\toprule
\textbf{Parámetro} & \textbf{Valor} \\
\midrule
Máximo de reintentos por modelo & 3 \\
Multiplicador de backoff & 2.0 \\
Espera mínima & 1\,s \\
Espera máxima & 60\,s \\
Timeout por llamada (Haiku) & 30\,s \\
Timeout por llamada (Sonnet) & 45\,s \\
Timeout por llamada (Opus) & 60\,s \\
\bottomrule
\end{tabularx}
\end{table}

El escalado se activa bajo tres condiciones: \texttt{RateLimitError} (error 429 de Anthropic), \texttt{APIError} (errores 5xx o timeout de conexión), y \texttt{ModelOverloadedError} (respuesta vacía o truncada).

% ====================================================================
\section{Sistema de prompts: estructura completa}
\label{sec:appendix_prompts}
% ====================================================================

El prompt del sistema del agente está estructurado en 10 secciones ordenadas que definen el comportamiento completo:

\begin{table}[H]
\centering
\caption{Estructura del prompt del sistema del agente unificado (10 secciones).}
\label{tab:appendix_prompt}
\footnotesize
\begin{tabularx}{\columnwidth}{clX}
\toprule
\textbf{N} & \textbf{Sección} & \textbf{Contenido} \\
\midrule
1 & IDENTIDAD & Nombre (ChatHCE), rol (asistente clínico), propósito y contexto de uso \\
2 & CONTEXTO & Dataset \ac{MIMIC}-IV-\ac{ED}: 222 pacientes, 6 tablas, urgencias hospitalarias \\
3 & HERRAMIENTAS & Documentación de las 3 herramientas con parámetros y ejemplos \\
4 & ANTI-ALUCINACIÓN & Prohibiciones absolutas, manejo de datos faltantes, citación, incertidumbre \\
5 & IDIOMA & Respuestas siempre en español, terminología médica apropiada \\
6 & SCHEMA BD & Esquema completo de las 6 tablas con nombres exactos de columnas \\
7 & HERR. DETALLADAS & Reglas \ac{SQL}, tipos de consulta, ejemplos de uso correcto e incorrecto \\
8 & FORMATO & Estructura de respuestas, uso de markdown, citación de fuentes \\
9 & GUÍAS CLÍNICAS & Valores de referencia normales, niveles de acuidad, interpretación \\
10 & SELECCIÓN HERR. & Cuándo usar cada herramienta, reglas obligatorias, docs indexados \ac{UCI} \\
\bottomrule
\end{tabularx}
\end{table}

La sección 10 (SELECCIÓN DE HERRAMIENTAS) fue la más crítica tras las mejoras al pipeline \ac{RAG}. Incluye la regla fundamental:

\begin{quote}
\footnotesize\itshape
``Si la pregunta NO contiene un subject\_id o stay\_id específico de paciente, SIEMPRE invoca \texttt{search\_clinical\_documents} antes de responder desde tu conocimiento general.''
\end{quote}

Y lista explícitamente los tres documentos indexados (Manual del Residente, Estándares \ac{UCI}, Libro de la \ac{UCI}) con 8 ejemplos de preguntas que deben activar la herramienta.

% ====================================================================
\section{Caché, rate limiting y performance monitor}
\label{sec:appendix_infra}
% ====================================================================

\subsection{CacheManager}

El \texttt{CacheManager} implementa caché en memoria con cinco tipos especializados, evicción \ac{LRU} y \ac{TTL} configurable. Es un Singleton compartido entre todos los componentes.

\begin{table}[H]
\centering
\caption{Tipos de caché implementados en el CacheManager.}
\label{tab:appendix_cache}
\footnotesize
\begin{tabularx}{\columnwidth}{llXl}
\toprule
\textbf{Tipo} & \textbf{TTL} & \textbf{Contenido} & \textbf{Evicción} \\
\midrule
\texttt{memory} & 300s & Respuestas generales del agente & \ac{LRU} \\
\texttt{llm\_responses} & 300s & Respuestas de Claude & \ac{LRU} \\
\texttt{embeddings} & 3600s & Embeddings de consultas frecuentes & \ac{LRU} \\
\texttt{query\_results} & 300s & Resultados de consultas a Supabase & \ac{LRU} \\
\texttt{ui\_data} & 60s & Datos de la interfaz Streamlit & \ac{LRU} \\
\bottomrule
\end{tabularx}
\end{table}

\begin{table}[H]
\centering
\caption{Configuración del CacheManager.}
\label{tab:appendix_cacheconfig}
\footnotesize
\begin{tabularx}{\columnwidth}{llX}
\toprule
\textbf{Parámetro} & \textbf{Valor} & \textbf{Descripción} \\
\midrule
\texttt{max\_cache\_size\_mb} & 512 MB & Tamaño máximo total \\
\texttt{cache\_ttl\_seconds} & 300s & \ac{TTL} por defecto \\
\texttt{cache\_cleanup\_interval} & 600s & Intervalo de limpieza \\
Estructura interna & \texttt{OrderedDict} & Evicción \ac{LRU} eficiente \\
Thread safety & \texttt{threading.RLock} & Un lock por tipo + lock global \\
Cálculo de tamaño & \texttt{pickle.dumps()} & Estimación en bytes \\
\bottomrule
\end{tabularx}
\end{table}

\subsection{Connection pool manager}

\begin{table}[H]
\centering
\caption{Configuración del pool de conexiones a Supabase.}
\label{tab:appendix_connpool}
\footnotesize
\begin{tabularx}{\columnwidth}{llX}
\toprule
\textbf{Parámetro} & \textbf{Valor} & \textbf{Descripción} \\
\midrule
\texttt{pool\_size} & 10 & Conexiones activas en el pool \\
\texttt{max\_overflow} & 20 & Conexiones adicionales bajo alta carga \\
\texttt{timeout} & 30s & Tiempo máximo de espera \\
\texttt{recycle} & 3600s & Tiempo de vida máximo de una conexión \\
Reintentos & 3 & Intentos de reconexión con backoff exponencial \\
\bottomrule
\end{tabularx}
\end{table}

\subsection{Performance monitor}

El decorador \texttt{@track\_performance} instrumenta automáticamente las operaciones críticas, registrando métricas en SQLite (\texttt{data/performance\_metrics.db}).

\begin{table}[H]
\centering
\caption{Métricas registradas por el Performance Monitor.}
\label{tab:appendix_perfmetrics}
\footnotesize
\begin{tabularx}{\columnwidth}{lX}
\toprule
\textbf{Métrica} & \textbf{Descripción} \\
\midrule
Tiempo de ejecución & Duración total en milisegundos \\
Tokens utilizados & Estimación de tokens consumidos por el \ac{LLM} \\
Modelo activo & Nombre del modelo Claude utilizado \\
Éxito/Error & Estado y tipo de error si aplica \\
Tipo de operación & \texttt{process\_message}, \texttt{database\_query}, \texttt{rag\_search} \\
Session ID & Identificador de sesión para trazabilidad \\
\bottomrule
\end{tabularx}
\end{table}

% ====================================================================
\section{Patrones de diseño implementados}
\label{sec:appendix_patrones}
% ====================================================================

\begin{table}[H]
\centering
\caption{Patrones de diseño implementados en ChatHCE con ubicación y propósito.}
\label{tab:appendix_patrones}
\footnotesize
\begin{tabularx}{\columnwidth}{llX}
\toprule
\textbf{Patrón} & \textbf{Componente} & \textbf{Propósito} \\
\midrule
Adapter & \texttt{ClaudeToolAdapter} & Adapta herramientas al formato tool-calling de Claude/LangChain \\
\addlinespace
Strategy & Búsqueda \ac{RAG} & Estrategias intercambiables: similarity, MMR, hybrid \\
\addlinespace
Singleton & \texttt{CacheManager} & Instancia única de caché compartida \\
\addlinespace
Singleton & \texttt{RateLimiter} & Estado compartido de rate limiting \\
\addlinespace
Singleton & \texttt{ImprovedRAGService} & Evita recarga del modelo CUDA en re-runs de Streamlit \\
\addlinespace
Factory & \texttt{create\_*\_tool()} & Funciones de fábrica para instanciar herramientas \\
\addlinespace
Decorator & \texttt{@retry\_on\_failure} & Reintentos con backoff exponencial (3 intentos, factor 2.0) \\
\addlinespace
Decorator & \texttt{@track\_performance} & Monitoreo automático de rendimiento \\
\addlinespace
Observer & \texttt{PerformanceMonitor} & Monitoreo de métricas y alertas de rendimiento \\
\addlinespace
Chain of Resp. & Fallback \ac{LLM} & Cadena Haiku 4.5 $\to$ Sonnet 4.5 $\to$ Opus 4 \\
\bottomrule
\end{tabularx}
\end{table}

% ====================================================================
\section{Stack tecnológico completo}
\label{sec:appendix_stack}
% ====================================================================

\begin{table}[H]
\centering
\caption{Stack tecnológico completo del sistema ChatHCE.}
\label{tab:appendix_stack}
\scriptsize
\begin{tabularx}{\columnwidth}{l>{\raggedright\arraybackslash}p{4.5cm}X}
\toprule
\textbf{Componente} & \textbf{Tecnología} & \textbf{Uso} \\
\midrule
Lenguaje & Python 3.11.14 & Lenguaje principal \\
Interfaz & Streamlit~\citep{streamlit2019} & Chat interactivo \\
Agente & LangChain 1.2.7~\citep{chase2022langchain} & Orquestación y herramientas \\
\ac{LLM} principal & Claude Haiku 4.5 & Agente conversacional \\
\ac{LLM} viz. & Claude Sonnet 4.5 & Generación de código \\
SDK Anthropic & anthropic 0.77.0 & Cliente \ac{API} de Claude \\
Base de datos & Supabase PostgreSQL 15~\citep{supabase2020,postgresql2024} & \ac{MIMIC}-IV-\ac{ED} + aplicación \\
Vector store & pgvector 0.8.0~\citep{pgvector2024} & Embeddings \ac{RAG} \\
Cliente BD & supabase-py & Acceso desde Python \\
Embeddings & sentence-transformers 5.2.2~\citep{reimers2019sentencebert} & all-MiniLM-L6-v2 (384d) \\
Evaluación & \ac{RAGAS} 0.4.3 & Evaluación pipeline \ac{RAG} \\
Visualización & Plotly~\citep{plotly2015} & Gráficas interactivas \\
Datos & pandas / numpy & Procesamiento de datos \\
PDF & Docling / PyPDF2 & Extracción de texto \\
DOCX & python-docx & Procesamiento Word \\
Validación & Pydantic~\citep{pydantic2023} & Validación de entradas \\
Logging & structlog & Logging estructurado \\
Testing & pytest & Framework de pruebas \\
\bottomrule
\end{tabularx}
\end{table}
